\section*{Bloque A. Fundamentos teóricos y conceptuales}

\noindent\textit{Preguntas 1 a 30. Evalúan la comprensión analítica de la tecnología de producción, las isocuantas, los rendimientos a escala, la minimización de costos, la dualidad, la oferta competitiva y los conceptos de beneficio y excedente del productor.}

\q{La función de producción $Q=f(K,L)$ representa:}
\begin{opts}
    \item el costo mínimo necesario para producir una cantidad de producto.
    \item el máximo producto factible con insumos dados y tecnología vigente.
    \item el beneficio económico máximo alcanzable en equilibrio de mercado.
    \item la cantidad mínima de insumos requerida para evitar pérdidas.
\end{opts}

\q{Una isocuanta es el conjunto de combinaciones de insumos que:}
\begin{opts}
    \item generan exactamente el mismo costo total de producción.
    \item reportan exactamente el mismo ingreso total de ventas.
    \item permiten alcanzar exactamente el mismo nivel de producto.
    \item producen exactamente el mismo beneficio económico neto.
\end{opts}

\q{La relación marginal de sustitución técnica (RMST) entre trabajo y capital se define como:}
\begin{opts}
    \item capital sustituible por una unidad más de trabajo manteniendo $Q$.
    \item variación total del producto ante cambios simultáneos en insumos.
    \item pendiente constante de la recta presupuestaria o recta isocosto.
    \item cociente directo entre el salario nominal y la renta del capital.
\end{opts}

\q{Si $PMg_L$ y $PMg_K$ son los productos marginales del trabajo y del capital, la RMST de trabajo por capital es:}
\begin{opts}
    \item $w/r$.
    \item $r/w$.
    \item $PMg_L/PMg_K$.
    \item $PMg_K/PMg_L$.
\end{opts}

\q{Una RMST decreciente a lo largo de una isocuanta implica que las isocuantas son:}
\begin{opts}
    \item líneas rectas.
    \item cóncavas respecto del origen.
    \item convexas respecto del origen.
    \item en forma de L.
\end{opts}

\q{Existen rendimientos constantes a escala cuando, al duplicar todos los insumos:}
\begin{opts}
    \item el producto aumenta menos del doble.
    \item el producto aumenta más del doble.
    \item el producto permanece constante.
    \item el producto se duplica exactamente.
\end{opts}

\q{Para la función Cobb-Douglas $Q=AK^\alpha L^\beta$, los rendimientos a escala están dados por:}
\begin{opts}
    \item $\alpha/\beta$.
    \item $\alpha+\beta$.
    \item $\alpha\beta$.
    \item $\max(\alpha,\beta)$.
\end{opts}

\q{Si en una función Cobb-Douglas $Q=AK^\alpha L^\beta$ se cumple $\alpha+\beta>1$, la tecnología presenta rendimientos:}
\begin{opts}
    \item decrecientes a escala.
    \item crecientes a escala.
    \item constantes a escala.
    \item negativos a escala.
\end{opts}

\q{La ley de los rendimientos marginales decrecientes establece que:}
\begin{opts}
    \item en el largo plazo todos los insumos productivos son variables.
    \item al duplicar ambos insumos el volumen de producción se duplica.
    \item la producción conjunta resulta siempre más barata que separada.
    \item con un factor fijo, el $PMg$ del factor variable acaba cayendo.
\end{opts}

\q{La relación entre producto medio ($PMe$) y producto marginal ($PMg$) implica que:}
\begin{opts}
    \item el $PMe$ alcanza su máximo cuando el $PMg$ es máximo.
    \item $PMg=PMe$ cuando el $PMe$ es cero.
    \item $PMg=PMe$ en el máximo del $PMe$.
    \item el $PMg$ es siempre mayor que el $PMe$.
\end{opts}

\q{La etapa II de la función de producción, considerada la región racional de producción, comprende:}
\begin{opts}
    \item desde el origen hasta el máximo del producto medio.
    \item desde el máximo del producto medio hasta que el producto marginal se hace cero.
    \item desde que el producto marginal se hace cero hasta que se vuelve negativo.
    \item toda la función de producción.
\end{opts}

\q{La recta isocosto $C=wL+rK$ tiene una pendiente igual a:}
\begin{opts}
    \item $-r/w$.
    \item $-w/r$.
    \item $-PMg_L/PMg_K$.
    \item $-P/Q$.
\end{opts}

\q{En el óptimo de minimización de costos para una producción dada, se cumple:}
\begin{opts}
    \item $RMST_{LK}=w/r$.
    \item $RMST_{LK}=r/w$.
    \item $PMg_L=PMg_K$.
    \item $wL=rK$.
\end{opts}

\q{La senda de expansión de la empresa es:}
\begin{opts}
    \item la ruta de combinaciones óptimas de insumos al variar $Q$ a precios fijos.
    \item la evolución del producto total al aumentar un único factor variable.
    \item la recta isocosto asociada exclusivamente a los costos fijos medios.
    \item la función inversa de demanda enfrentada por la firma competitiva.
\end{opts}

\q{La demanda condicionada de un factor productivo se obtiene al:}
\begin{opts}
    \item maximizar el beneficio económico dado un precio de mercado.
    \item minimizar el costo para un nivel de producto y precios dados.
    \item maximizar la producción total sujeta a un presupuesto fijo.
    \item determinar simultáneamente la curva de oferta del mercado.
\end{opts}

\q{El lema de Shephard aplicado a la función de costos establece que:}
\begin{opts}
    \item $\partial C/\partial Q=CMg$.
    \item $\partial C/\partial w=L(w,r,Q)$.
    \item $\partial C/\partial L=w$.
    \item $\partial Q/\partial L=w$.
\end{opts}

\q{Existen economías de escala cuando, al aumentar la producción:}
\begin{opts}
    \item el costo medio aumenta.
    \item el costo marginal aumenta necesariamente.
    \item el costo medio disminuye.
    \item el costo total disminuye.
\end{opts}

\q{Existen economías de alcance cuando:}
\begin{opts}
    \item la empresa fabrica exclusivamente un único bien homogéneo.
    \item el costo medio de un bien cae al elevar su propia escala.
    \item el costo marginal resulta constante en todo nivel de $Q$.
    \item producir dos bienes juntos cuesta menos que por separado.
\end{opts}

\q{Un costo hundido es aquel que:}
\begin{opts}
    \item puede recuperarse íntegramente si la firma decide cerrar.
    \item coincide exactamente con el costo marginal de corto plazo.
    \item ya ocurrió, es irrecuperable y no debe guiar decisiones futuras.
    \item mide el costo de oportunidad del capital financiero propio.
\end{opts}

\q{El costo económico de producción se diferencia del costo contable porque:}
\begin{opts}
    \item incluye única y exclusivamente los egresos contables explícitos.
    \item excluye sistemáticamente la retribución normal al capital propio.
    \item incluye tanto los costos explícitos como los de oportunidad.
    \item excluye la depreciación económica de activos fijos productivos.
\end{opts}

\q{El beneficio normal se considera:}
\begin{opts}
    \item parte del costo económico al remunerar el capital propio.
    \item una ganancia extraordinaria superior al costo de oportunidad.
    \item un costo fijo contable irrelevante para la permanencia.
    \item una renta económica pura derivada de barreras de entrada.
\end{opts}

\q{En el corto plazo, una empresa competitiva decide suspender la producción cuando:}
\begin{opts}
    \item $P<CVMe$.
    \item $P<CMe$.
    \item $P<CMg$.
    \item $P<CFMe$.
\end{opts}

\q{En el largo plazo, una empresa competitiva decide salir del mercado cuando:}
\begin{opts}
    \item $P<CVMe$.
    \item $P<CMe$.
    \item $P<CMg$.
    \item $P<CFMe$.
\end{opts}

\q{La condición de maximización de beneficios de una empresa perfectamente competitiva es:}
\begin{opts}
    \item $P=CMe$.
    \item $IMg=CMe$.
    \item $P=CMg$.
    \item $P=CVMe$.
\end{opts}

\q{La curva de oferta de una empresa competitiva en el corto plazo coincide con:}
\begin{opts}
    \item toda la curva de costo medio total en su tramo decreciente.
    \item el tramo del costo marginal sobre el mínimo del $CVMe$.
    \item toda la curva de costo variable medio a partir del origen.
    \item la curva inversa de demanda enfrentada en el mercado local.
\end{opts}

\q{En el equilibrio de largo plazo de una industria perfectamente competitiva, se cumple:}
\begin{opts}
    \item $P=CMg=CMe$ con beneficio económico nulo.
    \item $P>CMe$ con beneficio económico positivo.
    \item $P<CVMe$ con cierre inmediato de planta.
    \item $P=CMg>CMe$ con rentas extraordinarias.
\end{opts}

\q{El excedente del productor se define como:}
\begin{opts}
    \item el ingreso total menos el costo total.
    \item el beneficio contable.
    \item el excedente del consumidor.
    \item el ingreso total menos el costo variable.
\end{opts}

\q{La dualidad entre producción y costos implica que:}
\begin{opts}
    \item la función de costos resulta independiente de la tecnología.
    \item la función de producción solo queda definida en corto plazo.
    \item dicha relación solo se cumple en tecnologías Cobb-Douglas.
    \item costos y producción son representaciones duales de la tecnología.
\end{opts}

\q{Si una función de producción es homogénea de grado uno, el teorema de Euler establece que:}
\begin{opts}
    \item $Q=PMg_K\cdot K+PMg_L\cdot L$.
    \item $Q=PMg_K\cdot K-PMg_L\cdot L$.
    \item $Q=PMe_K\cdot K+PMe_L\cdot L$.
    \item $Q=w\cdot L+r\cdot K\cdot P$.
\end{opts}

\q{La elasticidad de sustitución entre factores se define como:}
\begin{opts}
    \item el cambio porcentual del producto dividido por el cambio porcentual de los insumos.
    \item el cambio porcentual de la relación $K/L$ dividido por el cambio porcentual de la RMST.
    \item el cambio porcentual del costo dividido por el cambio porcentual del producto.
    \item el cambio porcentual del precio dividido por el cambio porcentual de la cantidad.
\end{opts}
